\cleardoublepage % memastikan bab baru dimulai di halaman ganjil (kanan)
\chapter{PENDAHULUAN}
\label{chap:pendahuluan}
% --- Latar Belakang ---
\section{Latar Belakang}
\label{sec:latar_belakang}

Pergeseran paradigma pendidikan tinggi global saat ini bergerak dari
pendekatan pedagogis yang berpusat pada pengajar menuju pendekatan yang berpusat
pada pelajar melalui \textit{Personalized Learning} \autocite{Ryoo2021}. Transformasi
ini didorong oleh sistem \textit{Adaptive Assessment} (AA) berbasis kecerdasan
artifisial yang memiliki kemampuan untuk secara dinamis menyesuaikan tingkat kesulitan konten berdasarkan kemampuan peserta didik secara hampir \textit{real-time} \autocite{Halkiopoulos2024}. Sejalan dengan visi tersebut, pemerintah Indonesia melalui Strategi Nasional Kecerdasan
Artifisial (Stranas KA) 2020--2045 mendorong pengembangan sistem
penilaian cerdas yang adaptif \autocite{BPPT2020}. Secara teknis, berbagai algoritma
mulai dari \textit{Bayesian Knowledge Tracing} (BKT), \textit{Deep Knowledge Tracing}
(DKT), hingga sistem \textit{Elo rating} telah dikembangkan untuk mengestimasi
kemampuan pelajar \autocite{Minn2022, Vesin2022}. Implementasi teknik-teknik ini krusial untuk meninggalkan pendekatan kaku \textit{one size fits all}
dan menuju pengoptimalan pembelajaran individu di era digital.

Meskipun teknik-teknik yang disebutkan sebelumnya unggul dalam mengoptimalkan jalur
belajar individu, penerapan yang naif terhadap algoritma yang hanya berfokus
pada kemampuan individu membawa risiko tersembunyi yang serius. Dalam
literatur sistem rekomendasi, fenomena ini dikenal sebagai masalah
\textit{filter bubbles}, dalam fenomena ini, algoritma secara tidak sengaja
mengisolasi
pengguna dari keragaman informasi dan perspektif dengan hanya menyajikan konten
yang selaras dengan riwayat mereka \autocite{Wang2022}. Dalam konteks pendidikan,
risiko ini menghasilkan fenomena \textit{overpersonalization} (personalisasi
berlebih), yang menyebabkan sistem terlalu agresif dalam menyesuaikan konten dengan
kenyamanan individu. Hal ini justru mengasingkan pelajar dari manfaat pembelajaran komunitas
\autocite{Halkiopoulos2024}. Algoritma yang bekerja semata-mata berdasarkan
prinsip homofili (kemiripan) cenderung menciptakan \textit{echo chambers} \autocite{Tommasel2021}, menghalangi keberagaman perspektif yang penting bagi
pengembangan keterampilan pelajar secara keseluruhan.

Tinjauan literatur terkini mengonfirmasi bahwa ketergantungan eksklusif pada
algoritma personalisasi yang hanya mengoptimalkan parameter individu tanpa ruang
eksplorasi membawa dampak kognitif yang merugikan bagi pelajar. Bukti
eksperimental menunjukkan bahwa lingkungan yang terpersonalisasi secara agresif
memicu pengambilan sampel informasi yang sangat selektif sehingga membentuk
representasi kognitif yang terdistorsi \autocite{Bahg2025}. Masalah ini semakin
kritis karena pelajar sering menunjukkan rasa percaya diri berlebih pada generalisasi
pengetahuan yang keliru dan rentan terjebak dalam \textit{learning plateau}
\autocite{Bahg2025, Halkiopoulos2024}. Dalam konteks nasional, jika implementasi Stranas KA hanya mengadopsi teknologi adaptif konvensional, ancaman \textit{filter
bubbles} akademik ini berisiko melahirkan generasi pelajar yang terlalu percaya diri dengan representasi kognitif yang keliru.
Oleh karena itu, risiko \textit{overpersonalization} ini harus dipahami dan dicari alternatif
intervensi teknologi adaptif yang tidak mengasingkan pelajar dari manfaat esensial
keberagaman perspektif.

% --- Rumusan Masalah ---
\section{Rumusan Masalah}
\label{sec:rumusan_masalah}
Berdasarkan latar belakang tersebut, dapat diidentifikasi adanya kebutuhan untuk
mengeksplorasi pengembangan sistem asesmen yang selaras dengan Stranas KA yang
mempertimbangkan risiko \textit{overpersonalization}. Kesenjangan teknologi saat ini terletak pada terbatasnya
penelitian yang menguji apakah integrasi fitur kolaboratif dalam sistem asesmen
adaptif dapat meminimalkan efek \textit{overpersonalization} melalui penyisipan keberagaman perspektif.
Maka dari itu, permasalahan utama yang ingin diselesaikan adalah sebagai berikut.
\begin{enumerate}
  \item Bagaimana merancang arsitektur sistem asesmen adaptif yang mampu
  mengintegrasikan parameter kemahiran individual dengan dimensi kontribusi
  sosial secara terpadu?
  \item Bagaimana mengembangkan mekanisme mitigasi \textit{overpersonalization}?
  \item Bagaimana mengimplementasikan rancangan arsitektur dan mekanisme tersebut
  menjadi purwarupa sistem asesmen adaptif kolaboratif yang fungsional?
\end{enumerate}

% --- Tujuan ---
\section{Tujuan}
\label{sec:tujuan}
Sebagai upaya untuk menjawab permasalahan yang disebutkan sebelumnya, penelitian ini bertujuan untuk \textbf{mengembangkan purwarupa sistem asesmen adaptif kolaboratif yang mengintegrasikan
mekanisme intervensi proaktif}, yang harapannya dapat berkontribusi pada alternatif solusi
penggabungan metrik kemahiran individual dengan dimensi kontribusi sosial ke
dalam sebuah arsitektur asesmen yang terpadu. Selain itu, dalam penelitian ini juga diimplementasikan pengembangan metode mitigasi
\textit{overpersonalization} melalui deteksi stagnasi pembelajaran secara otomatis,
yang diharapkan dapat menjaga dinamika pertumbuhan pengetahuan pelajar dalam
lingkungan yang terpersonalisasi.

Secara spesifik, tujuan penelitian ini adalah sebagai berikut.
\begin{enumerate}
  \item Merancang arsitektur sistem asesmen adaptif yang mampu mengintegrasikan
  parameter kemahiran individual (\textit{Elo rating}) dengan dimensi kontribusi
  sosial secara holistik.
  \item Mengembangkan mekanisme mitigasi \textit{overpersonalization}.
  \item Mengimplementasikan rancangan arsitektur dan mekanisme tersebut menjadi
  purwarupa sistem asesmen adaptif kolaboratif yang fungsional sebagai pembuktian
  konsep (\textit{proof-of-concept}).
\end{enumerate}


% --- Batasan Masalah ---
\section{Batasan Masalah}
\label{sec:batasan_masalah}
Untuk memastikan penelitian ini tetap fokus dan dapat diselesaikan dalam jangka waktu yang ditentukan, ditetapkan batasan-batasan sebagai berikut.
\begin{enumerate}
  \item Permasalahan asesmen adaptif yang
  dikaji difokuskan pada konteks yang selaras dengan visi Strategi Nasional
  Kecerdasan Artifisial (Stranas KA) Indonesia. Ini berarti penelitian berfokus
  pada pengembangan sistem yang mendukung ekosistem kolaboratif, yang disesuaikan
  dengan kebutuhan dan potensi implementasi di lingkungan pendidikan tinggi di
  Indonesia.
  \item Fokus utama penelitian ini adalah pada kesenjangan integrasi
  algoritmik antara asesmen adaptif individual dan fitur kolaboratif. Penelitian
  tidak menginvestigasi masalah lain dalam asesmen adaptif (seperti deteksi
  gaya belajar atau deteksi kecurangan dalam asesmen), melainkan terbatas pada
  perancangan model yang menggabungkan asesmen adaptif individual dengan fitur
  kolaboratif.
  \item Permasalahan yang dikaji terbatas pada penilaian luaran kognitif
  (\textit{cognitive outcomes}), yaitu pelacakan keterampilan dan pengetahuan
  pelajar. Penelitian ini tidak mencakup perancangan model untuk menilai luaran
  afektif (seperti motivasi atau kepuasan) atau luaran perilaku (seperti durasi
  belajar) sebagai ukuran utama kemahiran.
  \item Fokus penelitian ini adalah pada integrasi asesmen adaptif dengan fitur kolaboratif untuk diversifikasi perspektif, sehingga akan diutamakan pengembangan dengan teknik dan alat yang efektif dan efisien secara sumber daya (baik sumber daya komputasi, waktu, uang, manusia, maupun data) yang utamanya untuk dapat membuktikan keberhasilan integrasi tersebut.
  \item Implementasi purwarupa sistem diujicobakan pada populasi sampel terbatas.
  Merujuk pada konteks pendidikan tinggi di Indonesia dan ketersediaan sumber daya,
  uji coba purwarupa difokuskan pada mahasiswa angkatan 2023 dari Program Studi Sistem dan Teknologi Informasi ITB.
\end{enumerate}

% --- Metodologi Pengerjaan TA ---
\section{Metodologi}

Pelaksanaan tugas akhir ini mengadopsi dan mengadaptasi metodologi yang digunakan oleh \textcite{Vesin2022} dalam penelitian mereka mengenai asesmen adaptif dan rekomendasi konten. Metodologi ini terdiri atas beberapa tahapan utama yang sistematis, dimulai dari studi pendahuluan hingga analisis data.

\begin{enumerate}
    \item \textbf{Studi Literatur:} Melakukan tinjauan sistematis terhadap penelitian terkait asesmen adaptif, pembelajaran kolaboratif, dan implementasi sistem serupa.
    
    \item \textbf{Perancangan dan Pengembangan Sistem:} Merancang arsitektur perangkat lunak serta mengimplementasikan modul-modul sistem, termasuk integrasi algoritma asesmen adaptif dan antarmuka pengguna kolaboratif.
    
    \item \textbf{Sesi Pengguna (Pengujian dan Evaluasi):} Melaksanakan studi dengan partisipan sebagai sampel populasi pengguna untuk mengumpulkan data kuantitatif dan kualitatif mengenai penggunaan dan pengalaman mereka terhadap purwarupa sistem. Data yang dikumpulkan dapat mencakup data interaksi (\textit{click-streams}, \textit{log files}), hasil asesmen (peringkat yang dihasilkan), serta umpan balik pengguna (misalnya melalui survei atau wawancara).
    
    \item \textbf{Analisis Data:} Menganalisis data yang terkumpul untuk mengevaluasi fungsionalitas purwarupa dan efektivitas sistem dalam konteks asesmen adaptif dan kolaboratif, guna menjawab rumusan masalah serta tujuan penelitian.
    
    \item \textbf{Penarikan Kesimpulan dan Pelaporan:} Menyusun laporan tugas akhir berdasarkan hasil analisis dan temuan penelitian.
\end{enumerate}

% --- Sistematika Penulisan ---
\section{Sistematika Penulisan}
\label{sec:sistematika_penulisan}
Sistematika penulisan ini memberikan gambaran umum mengenai struktur dan alur pembahasan dalam penyusunan laporan tugas akhir ini:
\begin{enumerate}
\item	\textbf{Bab I Pendahuluan:} Berisi latar belakang masalah, rumusan masalah, tujuan penelitian, batasan masalah, metodologi pelaksanaan, dan sistematika penulisan laporan tugas akhir.
\item	\textbf{Bab II Studi Literatur:} Berisi kajian pustaka yang relevan dengan topik tugas akhir, termasuk teori, konsep, dan hasil penelitian-penelitian terdahulu yang mendukung perancangan dan evaluasi sistem.
\item	\textbf{Bab III Analisis:} Berisi penjelasan mengenai analisis kebutuhan sistem asesmen adaptif dan kolaboratif, serta berbagai alternatif solusi untuk mitigasi \textit{overpersonalization}.
\item   \textbf{Bab IV Perancangan Arsitektur dan Interaksi Sistem:} Berisi penjelasan mengenai perancangan arsitektur sistem yang diusulkan, rancangan algoritma, desain basis data, dan antarmuka pengguna (\textit{user interface}).
\item	\textbf{Bab V Implementasi Purwarupa Sistem:} Berisi penjelasan mengenai proses implementasi solusi berdasarkan hasil perancangan, termasuk lingkungan pengembangan, batasan sistem, dan dokumentasi modul atau komponen yang dibangun.
\item	\textbf{Bab VI Evaluasi:} Berisi analisis dan evaluasi terhadap purwarupa sistem yang telah diimplementasikan, meliputi metodologi pengujian dengan pengguna, analisis data hasil pengujian kuantitatif dan kualitatif, serta validasi pencapaian tujuan.                  
\item   \textbf{Bab VII Penutup:} Berisi kesimpulan dari seluruh tahapan dan hasil pelaksanaan tugas akhir beserta saran yang dapat digunakan untuk pengembangan sistem lebih lanjut.
\item   \textbf{Daftar Pustaka:} Berisi daftar referensi yang menjadi rujukan dalam penyusunan laporan tugas akhir.
\item   \textbf{Lampiran:} Berisi dokumen pendukung penelitian seperti rincian \textit{source code} sistem dan data mentah hasil pengujian. 
\end{enumerate}
